\section{Discussion}

Our results build on the proximity framework established by Guney et al.\ \citep{Guney2016}: the proximity Z-score is a calibrated \textit{evidence} statistic---each compound is compared to a size- and degree-matched null, and a larger target set legitimately yields sharper significance. The phenomenon we characterize arises when the task shifts from per-compound classification to cross-compound ranking. Under target-count asymmetry, Z-score \textit{magnitude} should not be read as a biological effect-size ordering. Because the Z-score's denominator shrinks as $|T|^{-1/2}$ (a property we confirm for shortest-path, random-walk, and expression-weighted influence alike), the evidence ranking can diverge from, and even reverse, the raw topological effect ranking. This is the established distinction between effect size and statistical significance emphasized by the American Statistical Association \citep{Wasserstein2016}, here in a network-medicine setting.

The practical recommendation is to report both the evidence Z-score and an explicit effect size. We use per-target random-walk influence (perturbation efficiency), which by linearity of the walk equals the mean single-target influence on the disease module and is therefore comparable across compounds with unequal target counts. This metric complements the Z-score rather than replacing it: it does not itself escape variance shrinkage at the level of its own Z-score. Random-walk influence is also sensitive to node degree \citep{Cowen2017, Erten2011}, which is one reason we decompose it into direct-overlap and propagated components below.

Perturbation efficiency ranks Hyperforin above Quercetin, consistent with its xenobiotic-metabolism biology; Hyperforin is the PXR/CYP-inducing constituent mechanistically linked to DILI-relevant bioactivation \citep{Moore2000, Watkins2001, Chen2022}. This ordering holds across restart-probability variation, network threshold, and degree-binning. Decomposing the advantage, however, shows that it is largely \emph{direct} target--disease overlap: 62\% of Hyperforin's per-target influence comes from four targets that are themselves DILI-module genes. The \emph{propagated} component, isolated by leave-one-out, is a modest $1.5\times$ (about $1.2$--$1.9\times$ across alternative exclusions); its distribution overlaps size-matched Quercetin subsets, though it remains above a degree-matched background. Both components concentrate in the xenobiotic-metabolism axis that defines the mechanism and the overlap alike. Hyperforin's propagated advantage is above a degree-matched background (99.9th percentile, $3.3\times$ the mean) and above the mean of size-matched Quercetin subsets, while still overlapping their upper tail; the direct-overlap component is mechanistically interpretable but not independent evidence of propagated network positioning.

\subsection{Relationship to prior work}

Our findings build on Guney et al.'s proximity framework \citep{Guney2016} by identifying a cross-compound ranking regime that the original per-pair study design did not stress-test. Guney et al.\ evaluated proximity as a per-pair classifier across 238 drugs and 78 diseases (402 known drug--disease associations; mean 3.5 targets per drug), and demonstrated that the closest-distance relative proximity is essentially uncorrelated with target degree ($\rho=-0.01$); that is a statement about classification calibration within each drug's own size-matched null, not about the comparability of Z-\textit{magnitudes} across compounds whose target counts differ several-fold. Our task is different: cross-compound ranking under extreme target-count asymmetry (10 vs 62). In that regime the size-dependence of the null standard deviation, which the per-pair Z-score calibrates through its null mean and null variance, re-emerges when Z-magnitudes are compared between compounds. The resolution is not a new significance statistic but the separate reporting of effect size, for which per-target influence is a natural choice. Random walk with restart is conceptually adjacent to attenuated-distance and diffusion-style measures, but it is not identical to Guney's diffusion-kernel proximity ($d_k$), which downweights longer \emph{shortest paths} by an exponential penalty rather than propagating restart-anchored probability mass over the adjacency matrix; we therefore do not use Guney's $d_k$ result to argue for or against random-walk influence. The choice between proximity measures is orthogonal to our argument: a normalized influence measure can define an effect-size scale that should be reported \emph{alongside} the standardized evidence rather than in place of it when compounds of unequal size are compared. Degree bias in network propagation and prioritization is well documented \citep{Erten2011, Cowen2017, Barel2020}; the operating-regime benchmark (\S\ref{subsec:opregime}) complements node-degree control by explicitly calibrating the target-count axis, establishing that the $|T|^{-1/2}$ null-precision law is stable across the DILI module and size-matched pseudo-modules and that the resulting effect/evidence reversal is a \emph{located, conditional} regime — one that materializes when the smaller target set is above the 90th percentile of probe-pair margins.

\subsection{Limitations}

This study is a controlled two-compound audit designed to characterize the statistical behavior of proximity Z-scores under target-count asymmetry; it does not evaluate toxicological outcomes or predict DILI. Hyperforin is not an established intrinsic hepatotoxin \citep{LiverToxSJW}; the pair was selected for target-count contrast and mechanistic interpretability. Network influence is a measure of topological reach, not a toxicological outcome: it omits dose, exposure, pharmacokinetics, binding directionality (agonism vs antagonism), and reactivity. The two target sets also differ in provenance, not only size: Hyperforin's ten targets are literature-curated and mechanistically central (the PXR receptor with its induced enzymes and transporters), whereas Quercetin's sixty-two are a broad ChEMBL bioactivity net (Table~\ref{tab:s_evidence}). This evidentiary asymmetry could in principle contribute to Hyperforin's apparent advantage independently of topology; the direct--propagated decomposition and the DILI-module sensitivity (Table~\ref{tab:s_modulesens}) are designed to bound that confound (the propagated advantage survives both), but they do not eliminate it, and a fully curation-matched comparison would require harmonized target evidence for both compounds, which we note as a limitation. The ChEMBL-derived Quercetin target set may also include assay-context targets that differ from clinically relevant exposure targets. Generalization to compound libraries, for example testing whether perturbation efficiency carries DILI signal beyond target count and degree across DILIrank 2.0 with curated target sets, is a direction for future work. The operating-regime benchmark (\S\ref{subsec:opregime}) establishes that the $|T|^{-1/2}$ exponent is stable across the DILI module and size-matched pseudo-modules and that the reversal is a located, conditional regime; the \emph{rate} of reversal depends on the module's distance geometry and would differ in other networks. Its probes are degree-controlled random sets rather than curated drug-target sets, so it characterizes the statistic, not a drug population. Finally, interactome incompleteness and degree bias in random-walk propagation \citep{Cowen2017, Erten2011} bound the precision of all network-based statements made here.

\subsection{Conclusion}

Network-proximity Z-scores are valid evidence statistics whose magnitude should not be read as a cross-compound effect-size ranking under target-count asymmetry. Reporting raw effect size, standardized evidence, and per-target influence decomposed into direct-overlap and propagated components separately yields a more transparent comparative picture. A degree-controlled liver-network calibration shows that the null-precision law is stable across the DILI module and size-matched pseudo-modules, and that material rank reversal is a conditional regime requiring a smaller target set above the 90th percentile of probe-pair margins. In the worked example, the framework ranks the PXR/CYP-inducing constituent higher while exposing and separating out the direct target--disease overlap that an uncritical headline number would have conflated with propagated influence.
